\documentclass[11pt,a4paper]{article}

\usepackage{kotex}
\usepackage{geometry}
\usepackage{amsmath}
\usepackage{amssymb}
\usepackage{booktabs}
\usepackage{array}
\usepackage{longtable}
\usepackage{hyperref}
\usepackage{enumitem}

\geometry{margin=24mm}
\hypersetup{
  colorlinks=true,
  linkcolor=blue,
  citecolor=blue,
  urlcolor=blue
}
\setlist{nosep}
\emergencystretch=3em

\title{\textbf{EmoNet v3.1: 신경 활성 추적을 감정 상태 표현으로 보는 동역학 기반 모델}\\
\large 한국어 상세 논문 원고}
\author{Working Draft}
\date{2026년 5월 3일}

\begin{document}
\maketitle

\begin{abstract}
본 논문은 EmoNet v3.1에서 제안한 핵심 가설, 즉 감정이 단일 라벨이나 응답 생성을 위한 설명문이 아니라 자극이 신경망 내부를 통과하며 남긴 활성 흐름이라는 관점을 정리한다. 본 연구에서 감정 상태 표현은 사람이 직접 작성한 감정 메타데이터가 아니라, 입력 자극이 네트워크 안에서 만든 시간별 활성, 지배적 흐름, 가지 텐서, 그리고 그 압축 표현으로 정의된다. 외부 감정 라벨과 평가 문장은 감정 그 자체가 아니라 이 내부 흐름이 감정 구조를 담고 있는지를 검증하기 위한 기준으로만 사용된다.

연구의 목표는 좋은 상담 응답을 만드는 것이 아니다. 목표는 ``신경 활성 추적이 감정 상태 표현으로 기능하는가''를 검증하는 것이다. 이를 위해 v3.1은 네 가지 증거를 구축했다. 첫째, 감정 흐름이 너무 짧게 죽거나 전체 뉴런이 과도하게 켜지지 않는 안정적인 동역학을 만들었다. 둘째, 그 흐름 안에 감정가, 사회적 방향, 행동 경향, 평가 계열, 통제감 상태 같은 감정축 정보가 남는지를 확인했다. 셋째, 흐름의 특정 감정축 방향을 교란하면 생성 응답의 감정축 해석도 같은 방향으로 이동하는지 확인했다. 넷째, 흐름 안의 특정 감정축 단서를 강하게 중립화하면 원래 감정축 신호가 약해지는지 확인했다.

최종 설정인 후반부 적응형 활성 조절 모델은 n=80 확인 실험에서 감정 흐름 길이 50.475, 한 칸짜리 흐름 비율 0.000, 활성 밀도 0.709412, 감정 구조 분리도 0.238547, 균형 보정 감정 신호 0.136426을 보였다. 또한 10개 기본 사례를 사용한 통합 인과 확인 실험에서 전체 성공률 0.916667, 강한 정보 중립화 성공률 0.975000, 방향 교란 성공률 0.775000, 동일 응답 무효 비교 tie 비율 1.000000을 얻었다. 이 결과는 EmoNet v3.1의 신경 활성 추적이 감정 상태 표현 후보로서 구조적 신호와 초기 인과적 조작 가능성을 모두 가진다는 것을 보여준다.
\end{abstract}

\tableofcontents
\newpage

\section{서론}

감정을 계산 모델로 다루는 전통적 방식은 대체로 두 방향으로 나뉜다. 하나는 문장이나 상황을 입력으로 받아 ``분노'', ``슬픔'', ``불안'', ``기쁨'' 같은 범주를 출력하는 분류 방식이다. 다른 하나는 감정을 긍정--부정, 각성--진정 같은 낮은 차원의 연속값으로 나타내는 방식이다. 최근 생성형 언어 모델에서는 감정 정보가 응답 생성을 위한 문체 조건, 공감 조건, 안전 조건, 또는 프롬프트 설명으로 사용되는 경우가 많다.

그러나 이러한 접근은 감정을 너무 빨리 출력 형식으로 바꾸는 경향이 있다. 감정은 단순히 이름 붙일 수 있는 라벨이 아니다. 같은 ``분노''라도 상대에게 따지고 싶은 분노, 관계를 회복하고 싶은 분노, 자신을 방어하고 싶은 분노, 상황을 피하고 싶은 분노는 서로 다르다. 같은 ``슬픔''이라도 자기 비난, 무력감, 상실, 도움 요청, 체념은 다른 내부 구조와 다른 행동 준비 상태를 가진다.

EmoNet v3.1은 바로 이 지점에서 출발한다. 이 연구가 묻는 질문은 다음과 같다.

\begin{quote}
감정은 최종 라벨인가, 아니면 자극이 신경망 내부를 통과하며 형성한 활성 흐름인가?
\end{quote}

본 연구는 두 번째 관점을 택한다. 감정 상태는 텍스트 밖에 붙어 있는 설명문이 아니라, 자극이 네트워크 내부에서 만든 시간적 흐름, 가지 구조, 활성 패턴, 그리고 그 패턴의 압축 표현으로 보아야 한다는 것이 v3.1의 중심 가설이다.

\subsection{사용자가 만들고자 한 것}

이 프로젝트의 핵심 의도는 ``LLM이 더 그럴듯한 답변을 하게 만들자''가 아니다. 물론 EmoNet의 다른 버전에서는 내부 흐름을 응답 생성에 활용하는 방향이 있었다. 그러나 v3.1의 목표는 더 근본적이다. 사용자가 만들고자 한 것은 감정을 단순한 프롬프트 장식이나 감정 라벨로 다루지 않고, 자극이 내부 네트워크를 통과하며 만든 신경 활성 추적 자체를 감정 상태 표현으로 보는 연구선이다.

따라서 v3.1에서 중요한 것은 다음 질문들이다.

\begin{enumerate}
  \item 자극이 들어왔을 때 네트워크 내부에 충분히 긴 감정 흐름이 생기는가?
  \item 그 흐름이 너무 빨리 죽거나, 반대로 전체 뉴런을 거의 다 켜는 과활성이 되지는 않는가?
  \item 그 흐름 안에 감정가, 통제감, 사회적 방향, 행동 경향 같은 감정 구조가 남는가?
  \item 그 흐름의 특정 감정축을 바꾸면 응답의 감정축 해석도 따라 움직이는가?
  \item 그 흐름 안의 특정 감정축 단서를 제거하거나 중립화하면 원래 감정축 신호가 약해지는가?
\end{enumerate}

이 질문들이 모두 만족되어야만 ``trace가 감정 상태 표현 후보이다''라는 주장을 할 수 있다.

\subsection{본 논문의 기여}

본 논문은 다음 네 가지 기여를 정리한다.

\begin{enumerate}
  \item 감정 상태를 외부 라벨이 아니라 신경 활성 흐름으로 정의하는 EmoNet v3.1의 계산적 가설을 제시한다.
  \item 감정 흐름 붕괴와 과활성을 동시에 피하기 위한 후반부 적응형 활성 조절 방식을 구현하고 검증한다.
  \item 신경 활성 흐름 안에 감정축 구조가 남아 있는지 확인하기 위한 표현 지표를 제안하고 측정한다.
  \item 방향 교란, 강한 정보 중립화, 동일 응답 무효 비교를 포함한 축 전용 눈가림 인과 평가를 통해 trace 조작이 응답 수준 감정축 해석에 영향을 준다는 증거를 제시한다.
\end{enumerate}

\section{배경과 문제 정의}

\subsection{감정 라벨과 감정 흐름의 차이}

감정 라벨은 평가하기 쉽다. 예를 들어 어떤 문장에 ``분노''라는 라벨을 붙이면 정확도, 재현율, F1 점수 같은 지표로 모델을 평가할 수 있다. 그러나 라벨은 감정의 생성 과정을 거의 설명하지 않는다. 사람은 어떤 상황을 보고 즉시 라벨만 갖는 것이 아니라, 대상, 원인, 책임, 통제 가능성, 타인과의 관계, 행동 준비 상태를 함께 형성한다.

EmoNet v3.1은 이 과정을 내부 흐름으로 본다. 여기서 ``흐름''이란 시간에 따라 어떤 뉴런들이 켜지고, 어떤 경로가 유지되고, 어떤 가지가 지배적으로 살아남는지를 뜻한다. 이 흐름은 단일 숫자나 라벨보다 풍부하다. 본 논문에서는 이 흐름을 ``신경 활성 추적''이라고 부른다.

\subsection{핵심 용어}

본 논문에서는 가능하면 한국어 용어를 사용한다. 코드와 직접 연결되는 경우에만 괄호 안에 영어 이름을 병기한다.

\begin{longtable}{p{0.27\linewidth}p{0.24\linewidth}p{0.41\linewidth}}
\toprule
한국어 이름 & 코드 이름 & 의미 \\
\midrule
\endfirsthead
\toprule
한국어 이름 & 코드 이름 & 의미 \\
\midrule
\endhead
신경 활성 추적 & trace & 자극이 네트워크 내부를 통과하며 남긴 시간별 활성 기록 \\
감정 흐름 가지 & branch & 활성 경로 중 감정 흐름을 구성하는 가지 구조 \\
가지 텐서 & branch tensor & 시간별/경로별 가지 정보를 담은 배열 표현 \\
지배적 흐름 & dominant branch & 한 사례에서 가장 강하게 유지된 활성 경로 \\
감정 흐름 길이 & mean branch len & 흐름이 얼마나 길게 유지되었는지 \\
한 칸짜리 흐름 비율 & len1 ratio & 흐름이 거의 바로 죽어 한 칸만 남은 비율 \\
감정 흐름 붕괴 & branch collapse & 내부 경로가 충분히 형성되지 못하고 짧게 끝나는 현상 \\
활성 밀도 & activation density & 전체 뉴런 중 활성화된 뉴런의 비율 \\
과활성 & overactivation & 너무 많은 뉴런이 켜져 감정별 차이가 흐려지는 상태 \\
감정 구조 분리도 & separation & 서로 다른 감정축 집단이 trace 공간에서 분리되는 정도 \\
균형 보정 감정 신호 & balanced lift & 데이터 불균형을 보정한 이웃 일관성 향상도 \\
방향 교란 & perturbation & trace의 특정 감정축 값을 다른 방향으로 바꾸는 실험 \\
정보 제거 & ablation & trace의 특정 감정축 정보를 없애거나 약화시키는 실험 \\
강한 정보 중립화 & strong neutralization & 해당 축 field뿐 아니라 관련 문장 단서까지 중립화하는 제거 실험 \\
축 전용 눈가림 판정 & axis-only blind judge & 좋은 답변 여부가 아니라 지정 감정축만 보게 한 평가 \\
동일 응답 무효 비교 & null same response & 같은 응답을 A/B로 넣어 judge가 억지 선택을 하는지 보는 통제 실험 \\
\bottomrule
\end{longtable}

\subsection{검증해야 할 주장}

v3.1의 주장은 한 문장으로 요약하면 다음과 같다.

\begin{quote}
신경 활성 추적은 감정 상태 표현 후보이며, 감정축 구조를 보존하고 조작 가능하다.
\end{quote}

이 주장은 세 단계로 나뉜다.

\begin{enumerate}
  \item \textbf{동역학 증거}: 감정 흐름이 안정적으로 형성되어야 한다.
  \item \textbf{표현 증거}: 형성된 흐름 안에 감정축 구조가 남아 있어야 한다.
  \item \textbf{인과 조작 증거}: 흐름을 바꾸거나 중립화하면 응답의 감정축 해석도 예측 가능한 방향으로 바뀌어야 한다.
\end{enumerate}

이 세 단계 중 하나라도 빠지면 주장은 약해진다. 예를 들어 흐름이 안정적이지만 감정 구조가 없다면 단순 활성 패턴일 뿐이다. 감정 구조가 있지만 조작해도 응답이 변하지 않는다면 생성 수준에서 의미 있는 상태 표현이라고 말하기 어렵다. 반대로 응답이 변하더라도 내부 trace가 불안정하면 표현 자체가 신뢰하기 어렵다.

\section{시스템 구성}

\subsection{전체 파이프라인}

EmoNet v3.1의 전체 흐름은 다음과 같다.

\[
\text{입력 문장}
\rightarrow
\text{자극 벡터}
\rightarrow
\text{신경망 전파}
\rightarrow
\text{활성 추적}
\rightarrow
\text{가지 텐서}
\rightarrow
\text{감정 상태 표현 후보}.
\]

여기서 외부 감정 라벨은 감정 자체가 아니라 검증 기준이다. 예를 들어 ``통제감 상태''라는 라벨은 네트워크 내부에 통제감이 직접 문자열로 들어 있다는 뜻이 아니다. 이 라벨은 내부 trace가 통제감이 낮은 사례와 높은 사례를 구별할 수 있는지를 확인하기 위한 probe이다.

\subsection{신경망 동역학}

EmoNet은 각 tick마다 뉴런의 활성 후보를 계산한다. 후보 활성은 자극 입력, 이전 활성, 내부 정렬, 피로, 억제, 기억 계수 등의 영향을 받는다. v3.1에서 핵심 문제는 두 가지였다.

\begin{enumerate}
  \item 흐름이 너무 빨리 죽으면 감정 상태를 표현할 시간이 없다.
  \item 흐름을 억지로 살리면 거의 모든 뉴런이 켜져 감정별 특이성이 사라진다.
\end{enumerate}

이를 각각 감정 흐름 붕괴와 과활성이라고 부른다.

\subsection{후반부 적응형 활성 조절}

v3.1에서 최종적으로 사용한 방법은 후반부 적응형 활성 조절이다. 핵심 아이디어는 단순하다.

\begin{quote}
처음에는 감정 흐름이 생기도록 충분히 열어 두고, 일정 시간이 지난 뒤에만 활성 밀도를 조절한다.
\end{quote}

초반부터 강하게 억제하면 흐름이 생기지 않는다. 반대로 끝까지 열어 두면 과활성이 된다. 그래서 v3.1은 일정 tick 이후에만 활성 밀도가 목표 범위를 넘는지 확인하고, 필요하면 후보를 줄인다.

구현상 추가된 주요 설정은 다음과 같다.

\begin{longtable}{p{0.35\linewidth}p{0.55\linewidth}}
\toprule
코드 설정 & 한국어 의미 \\
\midrule
\endfirsthead
\toprule
코드 설정 & 한국어 의미 \\
\midrule
\endhead
\texttt{density\_control\_start\_tick} & 활성 밀도 조절을 시작하는 시점 \\
\texttt{density\_target\_high} & 이 값을 넘으면 부드럽게 후보를 줄이기 시작하는 목표 상한 \\
\texttt{density\_soft\_k\_leak\_gain} & 활성 후보를 부드럽게 낮추는 정도 \\
\texttt{density\_hard\_cap} & 절대 넘지 않도록 제한하는 활성 밀도 상한 \\
\texttt{density\_pruned\_fatigue\_gain} & 잘린 후보에 피로를 추가해 반복 과활성을 막는 정도 \\
\bottomrule
\end{longtable}

최종 best 설정의 코드 이름은 다음과 같다.

\begin{quote}
\texttt{adaptive\_thr0.63\_clip1.6\_inh0.10\_start8\_cap0.76}
\end{quote}

한국어로 풀면, 기본 발화 문턱은 0.63, 입력 신호 제한은 1.6, 억제 강도는 0.10, 8 tick 이후부터 활성 밀도 조절을 시작하고, 후반부 활성 밀도 상한을 0.76 근처로 제한하는 설정이다.

\section{평가 지표}

\subsection{동역학 지표}

동역학 지표는 trace가 감정 표현으로 쓰일 수 있을 만큼 건강하게 형성되는지를 본다.

\begin{enumerate}
  \item \textbf{감정 흐름 길이}: 흐름이 길수록 자극이 네트워크 안에서 충분히 전파되었다고 본다.
  \item \textbf{한 칸짜리 흐름 비율}: 이 값이 높으면 많은 사례에서 흐름이 거의 바로 죽었다는 뜻이다.
  \item \textbf{활성 밀도}: 너무 낮으면 흐름이 약하고, 너무 높으면 감정별 특이성이 사라진다.
\end{enumerate}

본 연구에서 활성 밀도의 작업 기준은 0.55에서 0.80 사이였다. 이 기준은 생물학적 상수가 아니라 실험적 기준이다. 너무 낮은 밀도에서는 흐름이 충분히 살아나지 않았고, 0.90 이상에서는 거의 전역 과활성처럼 보였기 때문이다.

\subsection{표현 지표}

표현 지표는 trace 안에 감정 구조가 남아 있는지 본다.

\begin{enumerate}
  \item \textbf{감정 구조 분리도}: 서로 다른 감정축 값들이 trace 공간에서 얼마나 떨어지는지 본다.
  \item \textbf{균형 보정 감정 신호}: 데이터에 많은 라벨이 이웃 찾기를 지배하지 않도록 class-balanced 방식으로 계산한 이웃 일관성 향상도이다.
\end{enumerate}

단순 이웃 일관성은 데이터 불균형에 약하다. 예를 들어 부정 감정 사례가 많으면, 아무 trace나 가까운 이웃을 찾아도 부정 라벨이 나올 가능성이 높다. 그래서 v3.1은 균형 보정 감정 신호를 추가했다. 이 값이 양수라는 것은 단순 다수 라벨 효과를 보정해도 trace 안에 감정축 구조가 남아 있다는 뜻이다.

\subsection{인과 지표}

인과 지표는 trace를 조작했을 때 응답 수준 감정축 해석이 바뀌는지를 본다.

\begin{enumerate}
  \item \textbf{방향 교란 성공률}: trace의 특정 감정축을 다른 값으로 바꾸었을 때, 응답도 그 방향으로 이동했는지 본다.
  \item \textbf{정보 중립화 성공률}: trace 안의 특정 감정축 단서를 중립화했을 때, 원래 감정축 신호가 약해지는지 본다.
  \item \textbf{동일 응답 무효 비교}: 같은 응답을 A/B로 넣었을 때 판정기가 tie를 내는지 본다.
\end{enumerate}

이 마지막 지표가 중요하다. 판정기가 같은 응답에서도 억지로 A나 B를 고른다면, 높은 성공률은 믿기 어렵다. 본 연구의 confirm10에서는 동일 응답 무효 비교가 40/40 tie로 나왔다.

\section{실험 과정}

\subsection{1단계: 기본 trace 추출}

첫 단계에서는 입력 사례를 네트워크에 넣고 시간별 활성, 지배적 흐름, 가지 텐서, 압축 표현을 저장했다. 이 단계의 목적은 응답 생성이 아니라 내부 흐름을 분석 가능한 형태로 꺼내는 것이었다.

추출된 핵심 산출물은 다음과 같다.

\begin{enumerate}
  \item 시간별 활성 기록
  \item 지배적 감정 흐름
  \item 가지 텐서
  \item 자극 벡터
  \item 활성 개수 통계
\end{enumerate}

\subsection{2단계: feature probe}

다음으로 어떤 trace feature가 감정 구조를 가장 잘 담는지 비교했다. 비교 대상은 단순 압축 벡터, 평균/최대 활성, 가지 평균, 시간 정보를 포함한 가지 표현 등이었다.

결론은 분명했다. 현재 가장 안정적인 감정 구조 신호는 단순 압축 벡터보다 가지 텐서 기반 표현, 특히 가지 평균에서 더 강하게 나타났다. 이는 v3.1의 핵심 직관과 잘 맞는다. 감정은 최종 압축 벡터 하나보다 자극이 통과한 경로와 가지 구조에 더 많이 남아 있었다.

\subsection{3단계: 용량 증가 실험}

뉴런 수를 늘리면 문제가 해결되는지 확인했다. 256개, 512개, 1024개 설정을 비교했다. 결과적으로 용량 증가는 감정 흐름 붕괴를 자동으로 해결하지 못했다. 더 큰 네트워크가 항상 더 좋은 trace를 만들지는 않았다.

이 결과는 중요하다. 문제는 단순히 모델 크기가 아니라 전파 동역학이었다. 흐름이 너무 빨리 죽거나 너무 많이 퍼지는 문제를 해결하지 않으면, 뉴런 수를 늘려도 감정 상태 표현은 안정되지 않는다.

\subsection{4단계: 동역학 안정화}

이후 문턱값, 피로, 억제, 입력 강도 등을 조정했다. 이 과정에서 중요한 tradeoff가 드러났다.

\begin{quote}
감정 흐름 붕괴를 없애면 활성 밀도가 너무 높아지고, 활성 밀도를 낮추면 흐름 붕괴가 다시 생긴다.
\end{quote}

예를 들어 지속 흐름 계열 설정은 한 칸짜리 흐름 비율을 0으로 만들 수 있었지만, 활성 밀도가 0.947까지 올라갔다. 이는 거의 대부분의 뉴런이 켜지는 상태이므로 감정별 특이성이 흐려질 위험이 컸다. 반대로 보수적 설정은 활성 밀도를 0.55 근처로 낮출 수 있었지만, 한 칸짜리 흐름 비율이 0.375--0.400 수준으로 남았다.

\subsection{5단계: 후반부 적응형 활성 조절}

이 tradeoff를 해결하기 위해 후반부 적응형 활성 조절을 추가했다. 이 방식은 초반 감정 흐름 형성은 보존하고, 후반 과활성만 조절한다.

n=40 adaptive sweep에서는 네 후보가 모두 한 칸짜리 흐름 비율 0.0과 목표 활성 밀도를 동시에 만족했다. 그중 가장 좋은 후보를 n=80으로 다시 확인했다.

\begin{table}[h]
\centering
\small
\begin{tabular}{p{0.34\linewidth}rrrrr}
\toprule
설정 & n & 흐름 길이 & 한 칸 비율 & 활성 밀도 & 감정 구조 분리도 \\
\midrule
\texttt{thr0.63 clip1.6 inh0.10 start8 cap0.76} & 80 & 50.475 & 0.000 & 0.709 & 0.239 \\
\bottomrule
\end{tabular}
\caption{최종 후반부 적응형 활성 조절 설정의 n=80 확인 결과.}
\label{tab:adaptive-confirm}
\end{table}

\subsection{6단계: 인과 평가 설계}

표현 증거만으로는 trace가 감정 상태 표현으로 기능한다고 충분히 말하기 어렵다. 그래서 trace를 조작했을 때 응답의 감정축 해석이 바뀌는지 보았다.

초기에는 trace 전체 응답과 정보 제거/교란 응답을 비교하는 방식으로 시작했다. 하지만 이때 LLM 판정기가 좋은 답변이나 자연스러운 답변을 고를 위험이 있었다. 이 문제를 줄이기 위해 축 전용 눈가림 판정을 만들었다.

축 전용 눈가림 판정은 다음 조건을 가진다.

\begin{enumerate}
  \item trace full, ablation, perturbation 같은 조건명을 숨긴다.
  \item A/B 순서를 섞는다.
  \item 친절함, 자연스러움, 유창함, 공감성, 전체 답변 품질을 평가하지 말라고 명시한다.
  \item 오직 지정된 감정축만 보게 한다.
  \item 동일 응답을 A/B로 넣어 무효 비교를 만든다.
\end{enumerate}

\subsection{7단계: 단순 정보 제거의 실패와 강한 정보 중립화}

초기 정보 제거는 약했다. 이유는 명확했다. 특정 field 하나를 지워도, preserve 문장이나 action tendency 문장 안에 같은 감정축 단서가 남아 있었다. 생성 모델은 그 남은 단서에서 원래 감정 방향을 복원할 수 있었다.

따라서 강한 정보 중립화를 만들었다. 이는 단순히 한 칸을 비우는 것이 아니라, 관련 문장 전체에서 해당 감정축 단서를 제거하거나 중립화한다.

예를 들어 통제감 상태를 중립화할 때는 통제감 field만 지우는 것이 아니라, ``무력감'', ``할 수 있음'', ``계획 가능성'', ``실행 단계'' 같은 통제감 단서가 preserve, avoid, action tendency 문장에 남지 않도록 바꾸었다.

\subsection{8단계: 통합 확인 실험}

마지막으로 같은 생성 조건에서 방향 교란과 강한 정보 중립화가 동시에 통과하는지 확인했다. 이 실험을 confirm10이라고 부른다.

설정은 다음과 같다.

\begin{enumerate}
  \item 기본 사례 10개
  \item trace full 10개
  \item 강한 정보 중립화 40개
  \item 일관 교란 40개
  \item 총 생성 응답 90개
  \item 축 전용 눈가림 판정 120개
  \item 생성기와 판정기: Claude Haiku 4.5
\end{enumerate}

여기서 생성기와 판정기가 같은 모델 계열이라는 점은 한계로 남는다. 그러나 같은 생성 조건에서 정보 중립화와 방향 교란이 동시에 기준을 넘었는지를 확인하는 목적에는 충분한 1차 확인 실험이다.

\section{결과}

\subsection{동역학 결과}

최종 설정은 감정 흐름 붕괴와 과활성을 동시에 해결했다.

\begin{table}[h]
\centering
\small
\begin{tabular}{lrr}
\toprule
지표 & 값 & 판정 \\
\midrule
감정 흐름 길이 & 50.475 & 좋음 \\
한 칸짜리 흐름 비율 & 0.000 & 붕괴 없음 \\
활성 밀도 & 0.709412 & 적정 범위 \\
감정 구조 분리도 & 0.238547 & 양수 \\
균형 보정 감정 신호 & 0.136426 & 양수 \\
\bottomrule
\end{tabular}
\caption{최종 설정의 핵심 동역학 및 표현 지표.}
\label{tab:core-results}
\end{table}

이 결과는 이전 tradeoff를 해결했다는 점에서 중요하다. baseline은 흐름이 짧았고, 지속 흐름 설정은 과활성이었으며, 보수적 설정은 흐름 붕괴가 남았다. 최종 adaptive 설정만 세 조건을 동시에 만족했다.

\subsection{감정축 표현 결과}

n=80 확인 실험에서 주요 감정축의 균형 보정 감정 신호는 모두 양수였다.

\begin{table}[h]
\centering
\small
\begin{tabular}{lr}
\toprule
감정축 & 균형 보정 감정 신호 \\
\midrule
감정가 & 0.258480 \\
사회적 방향 & 0.119748 \\
행동 경향 & 0.017921 \\
평가 계열 & 0.078728 \\
통제감 상태 & 0.175607 \\
\bottomrule
\end{tabular}
\caption{n=80 확인 실험의 축별 균형 보정 감정 신호.}
\label{tab:balanced-lift}
\end{table}

행동 경향은 다른 축보다 약하지만 양수이다. 이는 해당 축이 trace 안에 전혀 없다는 뜻이 아니라, 더 많은 데이터와 더 세밀한 행동 경향 분류가 필요하다는 뜻으로 해석하는 것이 안전하다.

\subsection{성능 비교}

주요 후보를 비교하면 다음과 같다.

\begin{table}[h]
\centering
\small
\begin{tabular}{lrrrr}
\toprule
설정 & 흐름 길이 & 한 칸 비율 & 활성 밀도 & 판정 \\
\midrule
baseline & 10.275 & 0.500 & 0.448 & 흐름 붕괴 \\
persistent less inhibition & 37.963 & 0.000 & 0.947 & 과활성 \\
conservative high fatigue & 19.825 & 0.400 & 0.550 & 붕괴 남음 \\
adaptive n=80 & 50.475 & 0.000 & 0.709 & 최종 선택 \\
\bottomrule
\end{tabular}
\caption{주요 동역학 설정 비교.}
\label{tab:dynamics-comparison}
\end{table}

최종 adaptive 설정은 baseline 대비 감정 흐름 길이를 크게 늘렸고, 한 칸짜리 흐름 비율을 0으로 낮췄으며, 활성 밀도를 목표 범위 안으로 넣었다.

\subsection{통합 인과 확인 결과}

confirm10 결과는 다음과 같다.

\begin{table}[h]
\centering
\small
\begin{tabular}{lrrr}
\toprule
평가 & n & 성공률 & tie 비율 \\
\midrule
전체 & 120 & 0.916667 & 0.358333 \\
강한 정보 중립화 & 40 & 0.975000 & 0.025000 \\
방향 교란 & 40 & 0.775000 & 0.050000 \\
동일 응답 무효 비교 & 40 & 1.000000 & 1.000000 \\
\bottomrule
\end{tabular}
\caption{confirm10 통합 인과 확인 실험.}
\label{tab:confirm10}
\end{table}

이 결과는 v3.1에서 가장 중요한 인과 증거이다. 같은 생성 조건에서 정보 중립화와 방향 교란이 모두 기준을 넘었다. 또한 동일 응답 무효 비교가 모두 tie였으므로, 판정기가 무조건 한쪽을 고르는 편향은 낮다.

축별 결과는 다음과 같다.

\begin{table}[h]
\centering
\small
\begin{tabular}{lrrr}
\toprule
감정축 & n & 성공률 & tie 비율 \\
\midrule
행동 경향 & 30 & 0.866667 & 0.400000 \\
통제감 상태 & 30 & 0.966667 & 0.333333 \\
사회적 방향 & 30 & 0.900000 & 0.366667 \\
감정이 향하는 대상 & 30 & 0.933333 & 0.333333 \\
\bottomrule
\end{tabular}
\caption{confirm10의 감정축별 결과.}
\label{tab:confirm10-axis}
\end{table}

\section{논의}

\subsection{무엇이 만들어졌는가}

v3.1에서 만들어진 것은 단순한 감정 분류기나 응답 생성 프롬프트가 아니다. 만들어진 것은 다음 세 가지를 포함하는 연구 시스템이다.

\begin{enumerate}
  \item 감정 상태를 신경 활성 추적으로 저장하고 분석하는 추출 체계
  \item 감정 흐름 붕괴와 과활성을 동시에 제어하는 후반부 적응형 활성 조절 방식
  \item trace의 감정축 구조와 조작 가능성을 검증하는 축 전용 눈가림 평가 체계
\end{enumerate}

이 세 요소가 결합되어야만 ``trace-as-emotion''이라는 주장을 할 수 있다. 단순히 trace를 저장하는 것만으로는 부족하고, 단순히 응답이 바뀌는 것만으로도 부족하다. 내부 흐름이 안정적이어야 하고, 그 흐름 안에 감정 구조가 있어야 하며, 그 흐름을 조작했을 때 응답 수준의 감정축 해석도 움직여야 한다.

\subsection{잘된 점}

첫째, 동역학 문제가 해결되었다. 감정 흐름이 짧게 죽는 문제와 과활성 문제가 동시에 해결되었다. 이는 v3.1의 가장 중요한 엔지니어링 성과이다.

둘째, 표현 증거가 확보되었다. 균형 보정 기준에서도 주요 감정축이 모두 양수였다. 이는 trace가 단순 노이즈가 아니라 감정 구조를 어느 정도 보존한다는 뜻이다.

셋째, 인과 조작 증거가 확보되었다. 방향 교란은 응답의 감정축 해석을 이동시켰고, 강한 정보 중립화는 원래 감정축 신호를 약화시켰다. 특히 동일 응답 무효 비교가 모두 tie였다는 점은 판정 결과의 신뢰도를 높인다.

\subsection{실패가 만든 개선}

v3.1은 한 번에 성공한 것이 아니다. 중요한 실패가 있었다.

초기 동역학 실험에서는 흐름 붕괴를 없애면 과활성이 생겼고, 과활성을 줄이면 흐름 붕괴가 돌아왔다. 이 실패 때문에 후반부 적응형 활성 조절이 필요해졌다.

초기 LLM 판정에서는 좋은 답변을 고르는 편향이 의심되었다. 이 실패 때문에 축 전용 눈가림 판정과 동일 응답 무효 비교가 추가되었다.

초기 정보 제거는 약했다. 특정 field 하나만 지워도 다른 문장에 같은 단서가 남아 있었기 때문이다. 이 실패 때문에 강한 정보 중립화가 도입되었다.

따라서 v3.1의 최종 구조는 단순히 성공한 실험들의 모음이 아니라, 실패를 통해 평가 기준과 조작 방식을 더 엄밀하게 만든 결과이다.

\subsection{주장 가능한 것과 아직 조심해야 할 것}

현재 결과로 주장할 수 있는 것은 다음이다.

\begin{quote}
EmoNet v3.1의 신경 활성 추적은 감정 상태 표현 후보로서 안정적인 동역학, 감정축 구조, 그리고 응답 수준 조작 가능성을 보인다.
\end{quote}

그러나 다음 주장은 아직 조심해야 한다.

\begin{enumerate}
  \item 모든 감정 현상에 대해 일반화된다고 말할 수는 없다.
  \item 사람 평가 없이 LLM 판정만으로 최종 심리학적 타당성을 주장할 수는 없다.
  \item 생성기와 판정기가 같은 모델 계열이므로 model-family bias 가능성이 남는다.
  \item full targeted set 전체나 외부 데이터셋에서 아직 확인하지 않았다.
\end{enumerate}

\section{한계}

본 연구의 한계는 명확하다.

\subsection{데이터 규모}

confirm10은 10개 기본 사례를 사용했다. dry3보다는 훨씬 강하지만, 대규모 검증이라고 부르기에는 아직 작다. 향후에는 전체 targeted set과 외부 데이터셋에서 확인해야 한다.

\subsection{판정기 독립성}

confirm10에서는 생성기와 판정기가 모두 Claude Haiku 4.5였다. 같은 모델 계열이 생성과 평가를 모두 담당하면, 특정 표현 방식에 유리한 편향이 생길 수 있다. 독립 LLM 판정기, 사람 평가, 또는 별도의 감정축 분류기를 추가해야 한다.

\subsection{감정축 정의}

target, social orientation, control state, action tendency class는 유용한 축이지만 최종적인 감정 이론 전체를 포괄하지 않는다. 특히 행동 경향은 더 세분화된 분류가 필요하다.

\subsection{응답 수준 평가의 한계}

응답은 trace뿐 아니라 생성 모델의 언어 습관, 안전성, 공손함, 상담 스타일에 영향을 받는다. 축 전용 눈가림 판정으로 이 문제를 줄였지만 완전히 제거한 것은 아니다.

\section{향후 연구}

향후 작업은 네 방향이다.

\begin{enumerate}
  \item \textbf{독립 판정}: 생성 모델과 다른 계열의 판정기 또는 사람 평가를 사용한다.
  \item \textbf{더 큰 표본}: confirm10을 full targeted set으로 확장한다.
  \item \textbf{직접 예측}: trace만 보고 감정축을 예측하는 별도 분류기를 학습한다.
  \item \textbf{세밀한 행동 경향}: 행동 경향 축을 더 촘촘하게 재정의한다.
\end{enumerate}

특히 가장 중요한 후속 작업은 다음 체인을 확인하는 것이다.

\[
\text{trace 조작}
\rightarrow
\text{trace 공간의 감정축 이동}
\rightarrow
\text{응답 수준 감정축 이동}.
\]

이 세 단계가 독립 판정기와 더 큰 데이터에서 재현되면, trace-as-emotion 주장은 훨씬 강해진다.

\section{결론}

EmoNet v3.1은 감정을 단일 라벨이나 응답 생성용 설명문이 아니라, 자극이 신경망 내부를 통과하며 만든 활성 흐름으로 보는 모델이다. 본 연구는 이 흐름이 감정 상태 표현 후보로 기능하기 위해 필요한 세 가지 조건, 즉 안정적인 동역학, 감정축 구조 보존, 인과적 조작 가능성을 검증했다.

최종 후반부 적응형 활성 조절 설정은 감정 흐름 붕괴와 과활성을 동시에 해결했다. n=80 확인 실험에서 한 칸짜리 흐름 비율은 0.000, 활성 밀도는 0.709412였고, 주요 감정축의 균형 보정 감정 신호는 모두 양수였다. 또한 confirm10 통합 인과 실험에서 강한 정보 중립화 성공률은 0.975000, 방향 교란 성공률은 0.775000, 동일 응답 무효 비교 tie 비율은 1.000000이었다.

따라서 v3.1의 결론은 다음과 같다.

\begin{quote}
신경 활성 추적은 감정 상태 표현 후보로서 구조적 신호를 가지며, 그 특정 감정축 단서를 교란하거나 중립화하면 응답 수준의 감정축 해석도 예측 가능한 방향으로 변한다.
\end{quote}

이 주장은 아직 인간 평가와 외부 데이터셋 검증이 필요한 1차 확정 결과이다. 그러나 v3.1은 이제 ``가능성 탐색'' 단계를 넘어, 논문화 가능한 연구 결과 단계에 도달했다.

\begin{thebibliography}{9}

\bibitem{russell2003core}
Russell, J. A. (2003).
Core affect and the psychological construction of emotion.
\textit{Psychological Review}, 110(1), 145--172.

\bibitem{lazarus1991emotion}
Lazarus, R. S. (1991).
\textit{Emotion and Adaptation}.
Oxford University Press.

\bibitem{scherer2001appraisal}
Scherer, K. R. (2001).
Appraisal considered as a process of multilevel sequential checking.
In \textit{Appraisal Processes in Emotion}.
Oxford University Press.

\bibitem{frijda1987action}
Frijda, N. H. (1987).
Emotion, cognitive structure, and action tendency.
\textit{Cognition and Emotion}, 1(2), 115--143.

\bibitem{picard1997affective}
Picard, R. W. (1997).
\textit{Affective Computing}.
MIT Press.

\bibitem{ledoux2012rethinking}
LeDoux, J. E. (2012).
Rethinking the emotional brain.
\textit{Neuron}, 73(4), 653--676.

\bibitem{barrett2017constructed}
Barrett, L. F. (2017).
\textit{How Emotions Are Made: The Secret Life of the Brain}.
Houghton Mifflin Harcourt.

\end{thebibliography}

\end{document}
